\begingroup
\setlength{\tabcolsep}{4pt}
\begin{tabularx}{\textwidth}{L{0.24\textwidth}X L{0.26\textwidth}}
\toprule
Tulajdonság & Jelentés & Tipikus példa \\
\midrule
Helyben rendezés & Az eredmény az eredeti tömbben áll elő, és az algoritmus csak konstans vagy kicsi extra tárat használ. & beszúró, gyorsrendezés, minimum kiválasztásos \\
Stabilitás & Az azonos kulcsú elemek megtartják egymáshoz képesti eredeti sorrendjüket. & beszúró, összefésülő, leszámláló \\
Adaptivitás & Az algoritmus kihasználja, ha a bemenet már részben rendezett. & beszúró, optimalizált buborék \\
Összehasonlító rendezés & A sorrend elempárok összehasonlításából derül ki. & beszúró, merge sort, quicksort \\
Nem összehasonlító rendezés & A kulcsok tartományát vagy szerkezetét is kihasználja. & leszámláló, számjegyes \\
Belső rendezés & Az adatok a főmemóriában vannak. & a legtöbb tömbös rendezés \\
Külső rendezés & Az adatok háttértáron vannak, ezért az I/O-költség is fontos. & külső összefésülés \\
\bottomrule
\end{tabularx}
\endgroup
